\section*{अध्याय \ref{ch:g10:coordgeom} --- निर्देशांक ज्यामिति}

\begin{solution}{exo:g10:coordgeom:1}
\omterm{prop:g10:coordgeom:midpoint}{मध्यबिंदु}: $\left(\frac{2 + (-4)}{2}, \frac{5 + 1}{2}\right) = (-1, 3)$।
दूरी:
\[
AB = \sqrt{(-4 - 2)^2 + (1 - 5)^2} = \sqrt{36 + 16} = \sqrt{52}
= 2\sqrt{13}.
\]
\end{solution}

\begin{solution}{exo:g10:coordgeom:2}
$I$ $[AB]$ का \omterm{prop:g10:coordgeom:midpoint}{मध्यबिंदु} है, इसका अर्थ है $\frac{3 + x_B}{2} = 1$ और $\frac{-2 + y_B}{2} = 2$, इसलिए $x_B = -1$ और $y_B = 6$: $B(-1, 6)$।
\end{solution}

\begin{solution}{exo:g10:coordgeom:3}
$OP = \sqrt{16 + 1} = \sqrt{17}$; $OQ = \sqrt{9 + 4} = \sqrt{13}$; $OR = \sqrt{0 + 25} = 5 = \sqrt{25}$। चूँकि $13 < 17 < 25$ है और वर्गमूल \omterm{def:g10:functions:variations}{वर्धमान} है, इसलिए मूल बिंदु के
सबसे निकट $Q$ है।
\end{solution}

\begin{solution}{exo:g10:coordgeom:4}
वर्ग-लंबाइयाँ:
\begin{align*}
AB^2 &= (5-1)^2 + (4-2)^2 = 16 + 4 = 20, \\
BC^2 &= (7-5)^2 + (0-4)^2 = 4 + 16 = 20, \\
AC^2 &= (7-1)^2 + (0-2)^2 = 36 + 4 = 40 .
\end{align*}
$AB = BC$: त्रिभुज $B$ पर समद्विबाहु है। इसके अतिरिक्त $AB^2 + BC^2 = 40 = AC^2$: पाइथागोरस प्रमेय
के विलोम से वह $B$ पर समकोण भी है।
\end{solution}

\begin{solution}{exo:g10:coordgeom:5}
\emph{1.} $[AC]$ का \omterm{prop:g10:coordgeom:midpoint}{मध्यबिंदु}: $\left(\frac{-2+4}{2}, \frac{1+1}{2}\right) = (1, 1)$; $[BD]$ का \omterm{prop:g10:coordgeom:midpoint}{मध्यबिंदु}: $\left(\frac{1+1}{2}, \frac{3+(-1)}{2}\right) = (1, 1)$। \omterm{prop:g10:coordgeom:midpoint}{मध्यबिंदु} एक ही
है: $ABCD$ समांतर चतुर्भुज है।

\emph{2.} $AB^2 = 3^2 + 2^2 = 13$ और $BC^2 = 3^2 + (-4)^2 = 25$: $AB \neq BC$, इसलिए समचतुर्भुज नहीं।

\emph{3.} $AC^2 = 6^2 + 0^2 = 36$ और $BD^2 = 0^2 + (-4)^2 = 16$: $AC \neq BD$, इसलिए आयत नहीं। $ABCD$ एक साधारण समांतर
चतुर्भुज है।
\end{solution}

\begin{solution}{exo:g10:coordgeom:6}
$\Omega(2,1)$ तक की वर्ग-दूरियाँ निकालिए और $5^2 = 25$ से तुलना कीजिए:
\[
\Omega A^2 = 3^2 + 4^2 = 25, \qquad
\Omega B^2 = (-3)^2 + (-4)^2 = 25, \qquad
\Omega C^2 = 4^2 + 1^2 = 17 .
\]
$A$ और $B$ वृत्त पर हैं; $C$ उसके भीतर है ($17 < 25$)।
\end{solution}

\begin{solution}{exo:g10:coordgeom:7}
$MA^2 = (x - 1)^2 + 1$ और $MB^2 = (x - 7)^2 + 25$। $MA^2 = MB^2$ रखने पर:
\[
x^2 - 2x + 2 = x^2 - 14x + 74
\ \Longleftrightarrow\
12x = 72
\ \Longleftrightarrow\
x = 6 .
\]
ऐसा अकेला बिंदु $M(6, 0)$ है।
\end{solution}

\begin{solution}{exo:g10:coordgeom:8}
\emph{1.} $I = \left(\frac{0+4}{2}, \frac{3+1}{2}\right) = (2, 2)$।

\emph{2.} $M(2,2)$ सचमुच $I$ है। और $MA^2 = (0-2)^2 + (3-2)^2 = 5$, $MB^2 = (4-2)^2 + (1-2)^2 = 5$: \omterm{prop:g10:coordgeom:midpoint}{मध्यबिंदु} दोनों सिरों से समान
दूरी पर है, जैसी अपेक्षा थी।

\emph{3.} ऐसा $N(0, y)$ खोजिए जिसके लिए $NA^2 = NB^2$ हो:
\[
(3 - y)^2 = 16 + (1 - y)^2
\ \Longleftrightarrow\
9 - 6y = 16 + 1 - 2y
\ \Longleftrightarrow\
-4y = 8,
\]
इसलिए $y = -2$: बिंदु $N(0, -2)$।
\end{solution}

\begin{solution}{exo:g10:coordgeom:9}
\begin{align*}
AB &= \sqrt{4^2 + 2^2} = \sqrt{20} = 2\sqrt5, \\
BC &= \sqrt{8^2 + 4^2} = \sqrt{80} = 4\sqrt5, \\
AC &= \sqrt{12^2 + 6^2} = \sqrt{180} = 6\sqrt5 .
\end{align*}
तब $AB + BC = 2\sqrt5 + 4\sqrt5 = 6\sqrt5 = AC$: त्रिभुज असमिका यहाँ समता बन जाती है, इसलिए $B$ रेखाखंड $[AC]$ पर है
--- तीनों बिंदु संरेख हैं।
\end{solution}

\begin{solution}{exo:g10:coordgeom:10}
$[AB]$ का \omterm{prop:g10:coordgeom:midpoint}{मध्यबिंदु} $I\left(\frac a2, \frac b2\right)$ है। तब
\[
OI^2 = \frac{a^2}{4} + \frac{b^2}{4},
\qquad
IA^2 = \left(a - \frac a2\right)^2 + \left(0 - \frac b2\right)^2
= \frac{a^2}{4} + \frac{b^2}{4},
\]
और $IB^2 = \left(0 - \frac a2\right)^2 + \left(b - \frac b2\right)^2 =
\frac{a^2}{4} + \frac{b^2}{4}$ भी। तीनों वर्ग-दूरियाँ बराबर हैं, इसलिए $OI = IA = IB$: समकोण त्रिभुज $OAB$ के
कर्ण का \omterm{prop:g10:coordgeom:midpoint}{मध्यबिंदु} तीनों शीर्षों से समान दूरी पर है (वही उसके परिवृत्त का
केंद्र है)।
\end{solution}

\begin{solution}{pb:g10:coordgeom:1}
\textbf{1.} दूरी सूत्र (\cref{thm:g10:coordgeom:distance}) से $\Omega M = 5 \iff \Omega M^2 = 25 \iff
(x - 2)^2 + (y - 1)^2 = 25$ --- वर्ग करना निरापद है, क्योंकि दोनों
पक्ष धनात्मक हैं।

\textbf{2.} $A$: $(5-2)^2 + (5-1)^2 = 9 + 16 = 25$: वृत्त पर। $B$: $16 + 9 = 25$: वह भी उस पर। $D$: $4 + 9 = 13 < 25$:
भीतर।

\textbf{3.} $x^2 - 6x = (x - 3)^2 - 9$ और $y^2 + 4y = (y + 2)^2 - 4$, इसलिए \omterm{def:g10:algebra:equation}{समीकरण} $(x - 3)^2 + (y + 2)^2 = 25$ बन जाता है: केंद्र $(3, -2)$,
त्रिज्या $5$।

\textbf{4.} $(x - 3)^2 + (y + 2)^2 = -1$: वर्गों का योग कभी ऋणात्मक नहीं होता --- कोई बिंदु इसे
संतुष्ट नहीं करता: रिक्त समुच्चय। सामान्य रूप से $(x - h)^2 + (y - k)^2 = c$ $c > 0$ होने पर त्रिज्या
$\sqrt c$ का वृत्त है, $c = 0$ होने पर अकेला बिंदु $(h, k)$, और $c < 0$ होने पर रिक्त।

\textbf{5.} $y = 3$ के साथ: $(x - 3)^2 + 25 = 25$, इसलिए $(x - 3)^2 = 0$: अकेला बिंदु $(3, 3)$। एक ही स्पर्श
बिंदु: रेखा वृत्त की \emph{स्पर्श रेखा} है।

\textbf{6.} $MA^2 + MB^2 = (x + r)^2 + y^2 + (x - r)^2 + y^2
= 2x^2 + 2y^2 + 2r^2 = 2r^2 + 2r^2 = 4r^2$ ($x^2 + y^2 = r^2$ का उपयोग करके)। चूँकि $AB^2 = (2r)^2 = 4r^2$ है, त्रिभुज $AMB$
$MA^2 + MB^2 = AB^2$ संतुष्ट करता है: वह $M$ पर समकोण है (पाइथागोरस का विलोम) --- तीन
पंक्तियाँ, जैसा वादा था।

\textbf{7.} $MA^2 + MB^2 = (x + a)^2 + y^2 + (x - a)^2 + y^2
= 2x^2 + 2y^2 + 2a^2 = 2\,MI^2 + 2a^2$, और $\frac{AB^2}{2} = \frac{(2a)^2}{2} = 2a^2$। सर्वसमिका सिद्ध।

\textbf{8.} $2\,MI^2 + \frac{36}{2} = 26$ से $MI^2 = 4$ मिलता है: केंद्र \omterm{prop:g10:coordgeom:midpoint}{मध्यबिंदु} $I$ और त्रिज्या
$2$ वाला वृत्त।

\textbf{9.} $MA^2 = 4\,MB^2$: $x^2 + y^2 = 4\left((x - 3)^2 + y^2\right)$, अर्थात् $3x^2 + 3y^2 - 24x + 36 = 0$, या $x^2 + y^2 - 8x + 12 = 0$: पूरा करने पर $(x - 4)^2 + y^2 = 4$ ---
केंद्र $(4, 0)$ और त्रिज्या $2$ वाला वृत्त (अनुपात $2$ का \emph{अपोलोनियस
वृत्त})।

\textbf{10.} यह पुल रचनाओं को गणनाओं में बदल देता है: न कोई सहायक आयत, न
स्थितियों का विश्लेषण --- एक प्रसार और प्रमेय हाथ में। क़ीमत यह है कि
गणनाएँ ज्यामितीय अंतर्दृष्टि कम देती हैं; चिरपरिचित उपपत्ति दिखाती है कि
\emph{क्यों}, और बीजगणितीय दिखाती है कि \emph{है} --- काम करने वाला गणितज्ञ
दोनों साथ रखता है।

\textbf{11.} $x^2 + y^2 = 25$ और $(x - 6)^2 + y^2 = 25$। घटाने पर: $x^2 - (x - 6)^2 = 0$, अर्थात् $12x - 36 = 0$: $x = 3$। फिर
$9 + y^2 = 25$: $y = \pm 4$। संभावित स्थितियाँ $(3, 4)$ और $(3, -4)$।

\textbf{12.} $(3, 4)$ से $R(0, 8)$ तक: $\sqrt{9 + 16} = 5$ --- मेल खाता है। $(3, -4)$ से: $\sqrt{9 + 144} \approx 12.4$ ---
ग़लत। आप $(3, 4)$ पर हैं।

\textbf{13.} दोनों \omterm{def:g10:algebra:equation}{समीकरणों} में वही $x^2 + y^2$ है (या प्रसार के बाद समान गुणांक
वाले पद): घटाने पर वर्ग कट जाते हैं और प्रथम घात का \omterm{def:g10:algebra:equation}{समीकरण} बचता है --- एक
रेखा। जब दोनों वृत्त दो बार मिलते हैं, तब वह रेखा दोनों मिलन-बिंदुओं से
होकर जाती है: वही उभयनिष्ठ जीवा की रेखा है।

\textbf{14.} चार अज्ञात: तीन \omterm{def:g10:coordgeom:system}{निर्देशांक} और अभिग्राही की घड़ी की त्रुटि
(सस्ती घड़ी, जिसे गणित सुधार देता है)। हर उपग्रह एक \omterm{def:g10:algebra:equation}{समीकरण} देता है, इसलिए
चार उपग्रह चार अज्ञातों के लिए चार \omterm{def:g10:algebra:equation}{समीकरण} देते हैं --- और अधिक उपग्रह
यथार्थता बढ़ा देते हैं।

\textbf{15.} अब घटाने पर $x = \frac{5.1^2 - 5^2 + 36}{12} = \frac{37.01}{12} \approx
3.084$ मिलता है: एक दूरी में $10$ cm की त्रुटि ने
यहाँ आकलन को लगभग $8$ cm खिसका दिया। त्रुटियाँ हर सूत्र में से होकर फैलती
हैं --- \cref{pb:g10:numbers:1} की अंतराल-अंकगणित वाली सावधानी, अब \omterm{def:g10:coordgeom:system}{निर्देशांकों} के साथ।

\textbf{16.} $B' (7, -3)$। रेखा $(AB')$ की \omterm{def:g10:reffunc:affine}{प्रवणता} $\frac{-3 - 5}{7 - 1} = -\frac43$ है और वह $y = 0$ को $x = 1 + \frac{5}{8} \times 6 = 4.75$ पर
काटती है: पानी $P(4.75,\ 0)$ पर पीजिए। \omterm{def:g10:functions:extrema}{न्यूनतम} लंबाई: $AP + PB = AP + PB' = AB' =
\sqrt{6^2 + 8^2} = 10$।

\textbf{17.} $Q(3, 0)$ होकर: $AQ + QB = \sqrt{4 + 25} + \sqrt{16 + 9} = \sqrt{29} + 5
\approx 10.39 > 10$। परावर्तन वाली सीधी रेखा जीतती है।

\textbf{18.} विकर्णों के \omterm{prop:g10:coordgeom:midpoint}{मध्यबिंदु}: $[AC]$: $(3, 3)$; $[BD]$: $(3, 3)$ --- बराबर:
समांतर चतुर्भुज। भुजाएँ: $AB = \sqrt{10} = AD$: समचतुर्भुज। विकर्ण: $AC = \sqrt{32} \neq BD = \sqrt{8}$: आयत नहीं, अतः
वर्ग भी नहीं। निर्णय: समचतुर्भुज, और इससे अधिक कुछ नहीं।

\textbf{19.} $D = A + C - B = (-1 + 6 - 3,\ 2 + 0 - 4) =
(2, -2)$: तब $[AC]$ और $[BD]$ का \omterm{prop:g10:coordgeom:midpoint}{मध्यबिंदु} एक ही $\left(\frac52, 1\right)$ है।

\textbf{20.} \omterm{prop:g10:coordgeom:midpoint}{मध्यबिंदु} सूत्र: समांतर चतुर्भुज और केंद्र (प्रश्न 18, 19)।
दूरी सूत्र: वृत्त, बिंदुपथ, त्रिभुजों की प्रकृति, सबसे छोटे रास्ते (प्रश्न
1--9, 16)। वृत्त-समीकरणों का घटाव: वह रेखा जो स्थिति बता देती है (प्रश्न
11--15)। बक्से में जो नहीं है: \emph{दिशाओं} का स्वच्छ बीजगणित --- सदिश ---
और \emph{प्रवणताओं} का --- रेखा-समीकरण तथा निकाय: अगले दो अध्याय ठीक यही
देते हैं।
\end{solution}
